\section{Méthode}
\label{sec:method}

\subsection{Formulation EWC et adaptation embarquée}
\label{sec:ewc_formulation}

\paragraph{Pénalité élastique.} EWC combat l'oubli catastrophique en ancrant les paramètres jugés
importants pour ce qui a déjà été appris \cite{Kirkpatrick2017EWC}. Après une tâche, on retient un
point d'ancrage $\theta^\star$ --- les paramètres tels qu'ils étaient à la fin de cette tâche --- et
une estimation $F$ de l'importance de chaque coordonnée. L'entraînement de la tâche suivante minimise
alors la perte de tâche augmentée d'un rappel quadratique vers cet ancrage :
\begin{equation}
\label{eq:ewc_loss}
\mathcal{L}(\theta) \;=\; \mathcal{L}_{\text{tâche}}(\theta)
\;+\; \frac{\lambda}{2} \sum_{i} F_i \,\bigl(\theta_i - \theta^\star_i\bigr)^2 .
\end{equation}
Le coefficient $\lambda$ arbitre entre plasticité et rétention. Chaque coordonnée est rappelée avec
une raideur qui lui est propre, $\lambda F_i$ : une coordonnée d'importance élevée devient coûteuse à
déplacer, une coordonnée d'importance faible reste libre. La somme porte sur \emph{tous} les
paramètres de la tête, poids et biais compris. Il n'y a qu'un seul terme quadratique, ancré sur le
dernier $\theta^\star$, et non une somme sur les tâches passées : le coût mémoire de la
régularisation est donc constant, deux copies des paramètres, quel que soit le nombre de tâches déjà
vues. Sur la première tâche, $F$ et $\theta^\star$ n'existent pas encore et la pénalité est absente.

\paragraph{Estimation de l'importance sur PC.} $F$ est la diagonale de la matrice d'information de
Fisher, estimée en fin de tâche sur les données de cette tâche. Nous en retenons la forme empirique,
c'est-à-dire évaluée sur les étiquettes réelles et non sur des étiquettes tirées sous le modèle. La
grandeur effectivement accumulée est le carré du gradient de la perte \emph{moyennée sur un lot},
moyenné sur les $B$ lots parcourus :
\begin{equation}
\label{eq:fisher}
F_i \;=\; \frac{1}{B} \sum_{b=1}^{B}
\left( \frac{\partial \mathcal{L}_b}{\partial \theta_i} \right)^{\!2},
\qquad
\mathcal{L}_b \;=\; \frac{1}{|b|} \sum_{n \in b} \ell\bigl(x_n, y_n; \theta\bigr) .
\end{equation}
Le budget est borné à \num{200} échantillons, ce qui suffit à ordonner les coordonnées entre elles ---
seul cet ordre compte, puisque $\lambda$ absorbe l'échelle globale. Deux écritures de l'accumulation
coexistent dans le projet. La variante \emph{en ligne} au sens de \cite{Schwarz2018Online} entretient
une importance unique en la faisant décroître, $F \leftarrow \gamma F + F^{(t)}$ avec $\gamma =
\num{0.9}$, ce qui estompe progressivement les tâches les plus anciennes. La tête effectivement
portée sur carte, elle, \emph{remplace} son importance à chaque consolidation : l'ancrage suit
toujours la dernière tâche consolidée.

\paragraph{Ce que la carte peut faire, et ce qu'elle ne peut pas.} Le portage conserve la pénalité
d'\eqref{eq:ewc_loss}, dont la dérivée $\lambda F_i (\theta_i - \theta^\star_i)$ s'ajoute au gradient
de tâche. La mise à jour embarquée s'écrit, pour un poids $W$ de la couche $l$ :
\begin{equation}
\label{eq:board_sgd}
W_{ji} \;\leftarrow\; W_{ji} \;-\; \eta \Bigl(
\underbrace{\delta_j \, a_i}_{\text{gradient de tâche}}
\;+\; \underbrace{\lambda \, F_{ji} \,\bigl(W_{ji} - W^\star_{ji}\bigr)}_{\text{rappel élastique}}
\Bigr),
\qquad \eta = \num{0.01} .
\end{equation}
La consolidation, en revanche, ne peut pas reproduire \eqref{eq:fisher} : l'estimer demanderait de
conserver un lot d'échantillons étiquetés et d'y repasser en rétropropagation, deux choses qu'une
carte en flux continu, sans mémoire de données, n'a pas à sa disposition. L'importance y est donc
approximée par la \emph{magnitude} du poids lui-même, entretenue en moyenne mobile exponentielle :
\begin{equation}
\label{eq:board_fisher}
F_{ji} \;\leftarrow\; \alpha \, F_{ji} \;+\; (1-\alpha) \, W_{ji}^2 ,
\qquad
W^\star_{ji} \;\leftarrow\; W_{ji} ,
\qquad \alpha = \num{0.99} .
\end{equation}
Il faut nommer cette substitution pour ce qu'elle est : $W^2$ n'est pas le Fisher, c'est un
substitut d'importance qui n'exige ni donnée, ni étiquette, ni passe arrière, et dont le coût est une
lecture-écriture par poids. Il fait l'hypothèse que les coordonnées de grande magnitude sont celles
qui portent la fonction apprise. Deux conséquences suivent. La pénalité embarquée ne couvre que les
matrices de poids, les biais n'ayant ni importance ni ancrage stockés, ce qui économise le tiers de
l'état de régularisation. Et à l'initialisation $F \equiv 0$ : tant qu'aucune consolidation n'a eu
lieu, le rappel élastique est nul et la carte apprend en descente de gradient pure.

\begin{table}[htbp]
\centering
\caption{Adaptation du modèle sous contrainte embarquée. La pénalité d'\eqref{eq:ewc_loss} est
conservée à l'identique ; ce sont son estimateur d'importance et son régime d'optimisation qui sont
simplifiés.}
\label{tab:ewc_port}
\footnotesize
\begin{tabular}{@{}p{30mm}p{35mm}p{50mm}@{}}
\toprule
 & \textbf{PC (PyTorch)} & \textbf{Carte (\texttt{ewc\_head.c})} \\
\midrule
Estimateur d'importance & carré du gradient, éq.~\eqref{eq:fisher} & magnitude $W^2$, éq.~\eqref{eq:board_fisher} \\
Données requises        & un lot étiqueté par tâche & aucune \\
Accumulation            & remplacement à chaque consolidation & moyenne mobile, $\alpha = \num{0.99}$ \\
Portée de la pénalité   & poids et biais & poids seulement \\
Optimisation            & lots, momentum, plusieurs époques & un échantillon, sans momentum, une passe \\
$\lambda$               & \num{1000} & \num{400} \\
Déclenchement           & fin de tâche & drapeau de trame, ou détection de dérive \\
État initial de $F$     & premier calcul en fin de tâche $1$ & nul, donc pas de rappel avant consolidation \\
\bottomrule
\end{tabular}
\end{table}

Un dernier écart, purement numérique, mérite d'être signalé car il éclaire les mesures de parité :
la carte met chaque poids à jour sur place puis propage l'erreur vers la couche précédente avec le
poids \emph{déjà modifié}, là où le cadre PC utilise partout les poids d'avant le pas. En inférence
gelée, les deux plateformes exécutent la même fonction et la parité est exacte ; en régime de mise à
jour, cet écart et le lot unitaire suffisent à expliquer qu'elle ne l'est plus tout à fait.

\subsection{Modèle EWC et portage C}

Le modèle est un perceptron multicouche à tête EWC (\texttt{EWCMlpMulticlass}, $k \to 32 \to 16 \to 2$)
entraîné en régime incrémental selon la Section~\ref{sec:ewc_formulation}. La tête est portée en C
(\texttt{ewc\_head.c}) pour la carte : le \emph{forward pass} et la mise à jour \gls{SGD} s'exécutent tous deux
sur le Cortex-M4 en FP32 (le F439ZI ne dispose pas de NPU). Les poids sont exportés du checkpoint PyTorch
vers un en-tête C généré (\texttt{model\_weights\_ewc.h}) --- jamais édité à la main --- de sorte que la
carte et le PC partagent \emph{exactement} les mêmes paramètres.

\subsection{Protocole apparié PC$\leftrightarrow$carte}

Pour rendre la comparaison honnête, la carte et le PC consomment la même séquence : même graine
($\texttt{seed}=42$), même ordre d'échantillons, même normalisation (les indices de features et les
tableaux sont produits par une source unique). Deux protocoles sont évalués : \emph{gelé} (inférence
seule, poids figés) et \emph{en ligne} (inférence puis mise à jour continue). La \emph{parité} est le taux
d'accord prédiction-à-prédiction entre les deux plateformes.

\subsection{Formulation de la quantification}
\label{sec:quant_formulation}

\paragraph{Le mapping et son échelle.} Quantifier un tenseur, c'est remplacer ses valeurs réelles par
des entiers sur une grille régulière, plus une échelle qui rend cette grille au monde flottant. Pour un
poids sur $b$ bits en représentation symétrique signée, avec $q_{\max} = 2^{\,b-1} - 1$ :
\begin{equation}
\label{eq:quant_map}
Q_{ji} = \mathrm{clip}\!\left( \mathrm{round}\!\left( \frac{W_{ji}}{s_j} \right),\, -q_{\max},\, q_{\max} \right),
\qquad
\widehat{W}_{ji} = s_j \, Q_{ji} .
\end{equation}
Tout le comportement du schéma tient dans le choix de $s_j$. Nous en comparons trois, qui forment une
seule famille et non des méthodes séparées :
\begin{equation}
\label{eq:scale_gran}
s_j \;=\;
\underbrace{\frac{1}{128}}_{\text{figée}}
\qquad\text{ou}\qquad
\underbrace{\frac{\max_{j,i} |W_{ji}|}{q_{\max}}}_{\text{par-tenseur}}
\qquad\text{ou}\qquad
\underbrace{\frac{\max_i |W_{ji}|}{q_{\max}}}_{\text{par-canal de sortie}} .
\end{equation}
L'échelle figée est celle du noyau embarqué d'origine. Elle ne dépend pas des poids : tout coefficient
de module supérieur à $127/128$ est écrasé sur la borne de la grille, et tous les coefficients plus
petits que le pas se confondent avec zéro.

\paragraph{La couche entière.} Les activations reçoivent le même traitement, avec une échelle $s_a$
issue d'une calibration. Une couche s'écrit alors comme une accumulation entière suivie d'un unique
retour au flottant :
\begin{equation}
\label{eq:int_layer}
\mathrm{acc}_j = \sum_i Q_{ji} \, q^a_i ,
\qquad
y_j = \mathrm{acc}_j \cdot s_j \cdot s_a + b_j .
\end{equation}
C'est dans \eqref{eq:int_layer} que se lit la granularité : l'indice $j$ porté par $s_j$ signifie qu'un
neurone de sortie dont les poids sont de faible amplitude reçoit sa propre échelle, au lieu de partager
celle du coefficient le plus grand de toute la matrice. La calibration retenue est un maximum absolu
par couche, mesuré sur un lot en précision flottante ($s_a = \mathrm{act\_max}/q^a_{\max}$).

\paragraph{Ce que le schéma d'origine prescrit, et ce que nous en gardons.} La formulation entière de
référence \cite{Jacob2018,Krishnamoorthi2018} vise une inférence \emph{entièrement} entière : le retour
à l'échelle suivante s'y fait par un multiplicateur entier et un décalage, les biais sont portés en
\texttt{int32} à l'échelle du produit, et un zéro-point décale la grille des activations. Notre cible
n'a pas les mêmes contraintes : le Cortex-M4F possède une unité flottante matérielle, et la tête ne
compte que trois couches. Nous conservons donc l'accumulation entière et les échelles par-canal, mais
nous revenons au flottant entre les couches. Le tableau~\ref{tab:quant_ref} récapitule les écarts.

\begin{table}[htbp]
\centering
\caption{Schéma entier de référence et mise en {\oe}uvre embarquée. Les écarts sont dictés par la cible
(unité flottante présente, réseau court) et non par une simplification de commodité.}
\label{tab:quant_ref}
\footnotesize
\begin{tabular}{@{}p{30mm}p{40mm}p{45mm}@{}}
\toprule
 & \textbf{Référence \cite{Jacob2018}} & \textbf{Noyau v2 (\texttt{ewc\_head\_int8\_v2.c})} \\
\midrule
Retour d'échelle    & multiplicateur entier et décalage & déquantification flottante sur l'unité matérielle \\
Biais               & \texttt{int32} à l'échelle du produit & flottants exacts, jamais quantifiés \\
Zéro-point          & sur les activations et les poids & poids symétriques sans zéro-point ; activations symétriques, affine en option \\
Granularité         & par-tenseur, par-canal en extension & par-canal de sortie par défaut \\
Calibration         & histogramme ou centile & maximum absolu par couche sur un lot \\
Fusion normalisation & requise avant quantification & sans objet, le réseau n'en contient pas \\
Accumulateur        & \texttt{int32} & \texttt{int32}, élargi en \texttt{int64} pour la variante Q15 \\
\bottomrule
\end{tabular}
\end{table}

\paragraph{Axe du format : l'échelle d'ablation.} Le noyau embarqué d'origine cumule quatre écarts à
\eqref{eq:quant_map}--\eqref{eq:int_layer} : échelle figée, arrondi remplacé par une troncature vers
zéro, accumulateur \texttt{int16} qui reboucle silencieusement, et activations repliées sur une grille
Q7. Plutôt que de les corriger d'un bloc, nous parcourons une trajectoire dans la famille
\eqref{eq:scale_gran}, chaque marche n'activant qu'un changement :
\begin{itemize}
  \item \texttt{legacy\_c} : échelle figée, troncature, accumulateur \texttt{int16} (le noyau v1
        d'origine sur carte) ;
  \item \texttt{fix\_acc32} : accumulateur élargi à \texttt{int32}, tout le reste inchangé ;
  \item \texttt{per\_tensor\_calib} : échelle par-tenseur \emph{calibrée} et arrondi au plus proche ;
  \item \texttt{per\_channel\_int8} : échelle \emph{par-canal} de sortie ;
  \item \texttt{q15} : grille 16 bits, $q_{\max} = 32767$ (RAM $\div 2$ au lieu de $\div 4$).
\end{itemize}
Le noyau v2 du firmware implémente les deux dernières marches, avec une variante Q15 optionnelle
(\texttt{-DEWC\_INT8\_Q15}) ; le chemin v1 reste inchangé pour permettre la comparaison sur la même
carte. L'ablation est \emph{émulée bit-à-bit sur PC} (\texttt{src/utils/int8\_c\_emulation.py}) afin
d'isoler la contribution de chaque correctif indépendamment de l'accès à la carte. Les échelles ne sont
pas recalculées à bord : elles sont produites sur PC par les mêmes fonctions que l'émulateur, puis
écrites dans un en-tête C généré, ce qui rend la parité vraie \emph{par construction} plutôt que
constatée après coup.

\paragraph{Où s'arrête la quantification en régime d'apprentissage.} Un point de portée doit être posé
avant de lire les cellules « en ligne » de la Section~\ref{sec:results}, car il en commande
l'interprétation. Le noyau entier n'expose qu'une passe avant : il n'existe pas de descente de gradient
dans le domaine quantifié. La carte conserve donc une tête FP32 \emph{maîtresse}, sur laquelle
s'appliquent \eqref{eq:board_sgd} et \eqref{eq:board_fisher}, et dont la représentation entière n'est
qu'une \emph{vue}, reconstruite par \eqref{eq:quant_map} après chaque pas. Deux conséquences en découlent,
et nous les assumons plutôt que de les taire. D'abord la division par quatre de la mémoire des poids vaut
pour l'inférence seule : en régime d'apprentissage, la carte porte simultanément la tête flottante, sa vue
entière, l'importance et l'ancrage, soit davantage que le FP32 seul. Ensuite la requantification de toutes
les matrices à chaque échantillon entre dans le budget de latence, ce que la
Section~\ref{sec:recovery-board} chiffre. Autrement dit, ce travail établit une quantification
\emph{compatible} avec l'apprentissage incrémental --- l'inférence est entière, l'apprentissage continue
de fonctionner à travers elle --- et non un apprentissage \emph{conduit} en arithmétique entière, qui reste
un problème ouvert.

\paragraph{Ce que coûte le par-canal.} Une échelle par neurone de sortie n'est pas gratuite : elle
ajoute un tableau de flottants parallèle à chaque matrice, soit une valeur sur quatre octets par ligne
de sortie. Sur la tête portée, ces échelles ramènent le gain mémoire réel des poids sous le facteur
quatre annoncé par le seul changement de type. C'est le prix de la robustesse : le par-canal isole les
neurones de grande amplitude, qui autrement imposeraient leur échelle à toute la couche.

\subsection{Axe du moment : quand quantifier}

À format fixé, la même tête peut être quantifiée à trois instants. \emph{Avant} l'entraînement, par QAT :
un \emph{fake-quant} inséré dans la boucle fait apprendre les poids sous la contrainte de précision
réduite. \emph{Après}, par PTQ : le modèle FP32 entraîné est converti une fois pour toutes. Et \emph{les
deux}, c'est-à-dire un entraînement QAT dont les poids sont ensuite exportés vers le kernel entier
calibré. Ce dernier chemin est le seul fidèle au déploiement réel, puisque c'est bien le noyau embarqué
qui exécute l'inférence, et non la simulation de quantification du cadre d'entraînement. Les trois
moments sont comparés à modèle, jeu et graine identiques.

\subsection{Axe de la profondeur : bits, granularité, symétrie}

Sous l'INT8, trois paramètres décrivent un schéma : le nombre de \emph{bits de poids} (4, 2, ternaire,
binaire), la \emph{granularité} de l'échelle (par-tenseur ou par-canal de sortie) et la \emph{symétrie}
(symétrique, ou affine avec zéro-point). Les deux premiers se lisent directement dans
\eqref{eq:quant_map} et \eqref{eq:scale_gran} : descendre en profondeur, c'est diminuer $q_{\max}$, et
changer de granularité, c'est changer la portée du maximum qui définit $s_j$.

Aux très basses profondeurs, la grille linéaire cède la place à deux schémas dédiés, que nous reprenons
tels quels. En ternaire \cite{Li2016TWN}, un seuil par canal décide quels poids survivent, et l'échelle
est la moyenne des seuls poids retenus :
\begin{equation}
\label{eq:twn}
\begin{gathered}
\Delta_j = \num{0.7} \cdot \overline{|W_{j\cdot}|},
\qquad
Q_{ji} = \mathrm{sign}(W_{ji}) \cdot \mathbf{1}\bigl[\,|W_{ji}| > \Delta_j\,\bigr], \\[2pt]
\alpha_j = \overline{\bigl\{ |W_{ji}| : |W_{ji}| > \Delta_j \bigr\}} .
\end{gathered}
\end{equation}
En binaire \cite{Rastegari2016XNOR}, plus aucun poids n'est mis à zéro et l'échelle est la moyenne des
modules du canal : $Q_{ji} = \mathrm{sign}(W_{ji})$, $\alpha_j = \overline{|W_{j\cdot}|}$. Ces deux
schémas portent une échelle par canal \emph{par construction}, si bien que l'axe de la granularité cesse
de s'appliquer une fois qu'on les atteint.

La symétrie, elle, concerne les activations. Après une fonction d'activation positive, la moitié
négative d'une grille symétrique ne sert à rien, ce qui coûte d'autant plus cher que les bits sont
rares. La variante affine décale la grille d'un zéro-point :
\begin{equation}
\label{eq:affine_act}
q = \mathrm{round}\!\left(\frac{a}{s_a}\right) + z,
\qquad
\widehat{a} = (q - z)\, s_a .
\end{equation}
Les poids, dont la distribution est centrée, restent symétriques sans zéro-point dans tous les schémas
que nous évaluons.

Un point mérite enfin d'être posé, car il commande
l'interprétation des gains annoncés : tant que les poids restent rangés dans un conteneur
\texttt{int8\_t}, réduire le nombre de bits ne libère \emph{aucun} octet. Le gain mémoire n'existe qu'avec
un \emph{bit-packing} effectif, qui se paie en revanche par un dépacking à l'exécution. Nous distinguons
donc systématiquement la RAM \emph{théorique} du schéma et la RAM \emph{mesurée} du binaire.

\subsection{Méthodologie de mesure}

Les grandeurs rapportées en Section~\ref{sec:system} appellent trois instrumentations distinctes.

\paragraph{Mémoire.} La RAM d'un binaire n'est pas l'empreinte de ses poids. Nous rapportons la
\emph{RAM totale} $= \texttt{.data} + \texttt{.bss} + \text{pic de pile}$, les deux premiers termes étant
lus dans le binaire et le troisième mesuré à l'exécution par marquage de la pile, en retenant le maximum
entre la phase d'inférence et la phase de mise à jour.

\paragraph{Latence.} Le compteur de cycles \gls{DWT} chronomètre le noyau au cycle près. Pour le chemin INT8, il
est instrumenté en \emph{trois segments} --- déquantification, \gls{MAC} entier, requantification --- afin de
répondre non pas « l'INT8 est-il plus lent ? » mais « quel poste coûte ? ». Les cycles sont rapportés
bruts, la conversion en microsecondes tronquant l'information à cette échelle.

\paragraph{Énergie.} Le courant est mesuré à la sonde X-NUCLEO-LPM01A par trois méthodes indépendantes :
régression du courant en fonction de la cadence d'inférence, différence contre un repos \gls{WFI}, et
lot d'inférences par trame isolant le calcul du coût de la liaison. Nous les rapportons \emph{côte à côte}
sans jamais les moyenner : elles ne mesurent pas la même frontière, et leur désaccord est lui-même un
résultat. Toute cellule non mesurée est reportée « à mesurer » plutôt qu'estimée.
